% ===================================================================
%  HOTO Na 4 — Balaga da tsufa.
%
%  Traduction, faite en 2026, en haoussa d'aujourd'hui, sur l'ido de
%  texto/io. Regles de la colonne : voir 10-hoto-01.tex. Deux scenes :
%  t04-c1-* la noce, t04-c2-* l'anniversaire du grand-pere.
%
%  « kaka » EST L'AUTOMNE ET C'EST AUSSI LE GRAND-PERE, ET CE TABLEAU
%  PORTE LES DEUX. Le meme mot, sans nuance ni derivation, nomme la
%  saison de la moisson et l'aieul : « kaka » l'automne a l'alinea 2,
%  « kaka » la grand-mere (29) a l'alinea 14, « kaka » le grand-pere
%  (33) a la seconde scene, et « kakan kaka » l'arriere-grand-pere au
%  dernier alinea — le grand-pere du grand-pere. QUATRE FOIS DANS UN
%  TABLEAU, et la planche a mis la saison et les vieillards sur la
%  meme feuille sans le savoir. Une langue ou l'annee et la famille
%  se comptent avec le meme mot n'a pas besoin de dire que l'automne
%  est le grand age.
%
%  ET LES CHRYSANTHEMES (9) SUIVENT : ils n'ont pas de nom en
%  haoussa, et la colonne les appelle « furannin kaka », les fleurs
%  de la moisson. C'est vrai de la fleur, c'est vrai de la saison, et
%  c'est vrai du repas de famille auquel elles sont posees.
%
%  L'HIVER EST LA SEULE SAISON QUI AIT DU S'ETIRER, ET C'EST L'EXACT
%  CONTRAIRE DU PERSAN. L'annonce du tableau 3 se tient : l'automne
%  tombe juste sur kaka, la moisson, puisque la planche grave des
%  vendanges. Mais l'hiver de la seconde scene est « hunturu », la
%  saison de l'harmattan — froide, seche, poussiereuse, et JAMAIS
%  BLANCHE. La neige (19) doit donc se construire : « dusar
%  ƙanƙara », la poudre de glace. La colonne persane, au meme tableau
%  et devant la meme planche, avait برف, یخ, زمستان, سرما — quatre
%  mots iraniens ordinaires, parce que Teheran neige tous les ans.
%  Meme planche, meme saison, deux langues, et l'une des deux ne l'a
%  jamais vue.
%
%  LE REMEDE EST HAOUSSA ET LE MEDECIN NE L'EST PAS. Tout ce qui se
%  soigne depuis toujours porte un nom d'ici : magani le remede, boka
%  le guerisseur, amosani le rhumatisme, zazzaɓi la fievre, mura le
%  rhume, tari la toux, jiri le vertige, suma l'evanouissement. Mais
%  « likita » le medecin (25) et sa « takardar magani » (27) sont
%  venus avec l'hopital. Le persan du meme tableau partageait sa
%  medecine entre le persan et l'arabe, et la ligne passait entre la
%  peau et les organes ; ici elle passe entre l'ancienne medecine et
%  la nouvelle. Deux traces differentes, une seule loi : ce qui s'est
%  ECRIT ailleurs arrive avec ses mots.
%
%  LA FOURCHETTE (12) EST « cokali mai yatsa », LA CUILLER A DOIGTS.
%  Le haoussa avait la cuiller — cokali — et n'avait pas la
%  fourchette ; il ne l'a pas empruntee, il l'a decrite par ce qui la
%  distingue. L'assiette (47), elle, est « faranti », qui est
%  « plate » refait a la bouche haoussa, le p devenu f comme partout.
%  La table de ce tableau est donc moitie decrite et moitie empruntee,
%  et le partage suit exactement ce que la maison avait deja.
%
%  DEUXIEME REFUS DE LA COLONNE : LE PARDESSUS (7). Le haoussa a
%  « babbar riga », la grande robe — un vetement precis, a coupe
%  precise, qu'un homme porte par-dessus tout le reste, et le mot
%  serait tombe tout seul sous la main. La planche grave un pardessus
%  d'hiver europeen a col de fourrure. La colonne ecrit « kwat na
%  sanyi », le manteau de froid. On modernise la langue, jamais les
%  choses — apres le bori du tableau 3, c'est le second.
%
%  LES PERDRIX (2) SONT « makwarwa », ET C'EST L'INVERSE DE L'OURS.
%  Le francolin de brousse est bien la, il vole en compagnie devant
%  le chien, et le haoussa a son nom depuis toujours. Quand la bete
%  est ici on la nomme ; quand elle n'y est pas — l'ours du tableau 2,
%  le tigre du tableau 3 — on emprunte et l'on n'en cache rien. Le
%  releve n'a pas encore vu de colonne qui invente une bete.
%
%  ET LE VIN DE LA NOCE SE DIT AVEC UN MOT PLUS VIEUX QUE L'ISLAM :
%  « giya », qui est la biere de mil et non un emprunt. Le champagne
%  (49), lui, n'a pas de nom et se decrit : « ruwan inabi mai
%  kumfa », l'eau de raisin qui mousse. Une langue musulmane depuis
%  mille ans qui a garde le nom de sa boisson d'avant.
% ===================================================================

%%K t04-apar-1 apar
\VUsaut{4.0mm}
\VUtitre{15.0pt}{60}{HOTO Na 4}
\VUsaut{1.6mm}
\VUtitre{11.4pt}{0}{Balaga da tsufa.}
\VUsaut{3.2mm}

%%K t04-tit-1 sub
\VUcentre{11.4pt}{0}{\textit{Fage na Farko.}}
\VUsaut{1.6mm}
\VUtitre{11.4pt}{0}{Bikin Aure.}
\VUsaut{2.6mm}

%%K t04-tit-2 sub
\VUpk{10.2pt}{40}{Kaka.}
\VUsaut{3.0mm}

%%K t04-c1-01-1 p
1. --- Samarin sun girma. Ɗaya daga cikinsu, Ioannes, yana yin aure.
Muna halartar \VUgras{bikin aure} wanda ya tara iyalan ango da amarya
kewaye da tebur ɗaya.

%%K t04-c1-02-1 p
2. --- Muna cikin \VUgras{kaka,} a \VUgras{watan Satumba.} Iskar sanyi
ta \VUgras{Oktoba} da ta \VUgras{Nuwamba} ba ta riga ta tuɓe ƙauyen
koren ganyensa ba. Iskar yamma tana da ɗumi da ƙamshi ; shi ya sa ake
cin abincin dare a ƙarƙashin \VUgras{baranda} \textsuperscript{(8)} kusa
da gona.

%%K t04-c1-03-1 p
3. --- Can nesa, a kan \VUgras{tudu} \textsuperscript{(3)}, akwai gidan
iyayen Ioannes, \VUgras{kagara} \textsuperscript{(1)} mai kwarjini ɓoye
cikin ganyaye ; kyawawan gonakin inabi, waɗanda suke ba da giya mai
kyau ƙwarai, sun rufe gefen tudun. Manomin inabi yana nomansu yana
kuma kula da su.

%%K t04-c1-04-1 p
4. --- A saman gonakin inabi, garken \VUgras{makwarwa}
\textsuperscript{(2)} yana wucewa, tsorace saboda kusatowar \VUgras{mai
gadin farauta} \textsuperscript{(5)}. Shi kuwa, da \VUgras{bindiga} a
hammata da \VUgras{jakar farauta} a kafaɗa, yana bincika
\VUgras{sarƙaƙƙiya} \textsuperscript{(4)} tare da \VUgras{karensa}
\textsuperscript{(6)}. Yana tsare gonar yana kuma hana ƴan farautar
sata su hallaka namun daji.

%%K t04-c1-05-1 p
5. --- A wajen \VUgras{baranda,} akwai \VUgras{shingen inabi} yana
hawa, wanda \VUgras{gungun inabi} \textsuperscript{(7)} masu yawa suke
rataye a kansa.

%%K t04-c1-06-1 p
6. --- Kyawawan furannin kaka, waɗanda suka yi wa \VUgras{tebur}
\textsuperscript{(16)} ado, su ne \VUgras{furannin kaka}
\textsuperscript{(9)} ; \VUgras{mahaifin} \textsuperscript{(10)} amaryar
ya sa ɗaya daga cikinsu a ramin maɓallin rigarsa.

%%K t04-c1-07-1 p
7. --- Cin abinci yana gab da ƙarewa. \VUgras{Mai hidima}
\textsuperscript{(11)} ya fara kawo kwanonin kayan zaki, ba da daɗewa
ba kuma zai ɗauke sassan kayan cin abinci iri iri, \VUgras{cokula}
\textsuperscript{(14)}, \VUgras{cokali mai yatsa} \textsuperscript{(12)},
\VUgras{wuƙaƙe} \textsuperscript{(13)}, \VUgras{manyan faranti}
\textsuperscript{(15)}, waɗanda ba a buƙatarsu kuma.

%%K t04-tit-3 sub
\VUpk{10.2pt}{40}{Dangi.}
\VUsaut{3.0mm}

%%K t04-c1-08-1 p
8. --- Kowa yana da farin ciki a fuska, murmushin da \VUgras{uwar}
\textsuperscript{(17)} angon take yi tana kallon ƴaƴanta kuwa yana
tabbatar da farin cikinta.

%%K t04-c1-09-1 p
9. --- Wannan aure yana haɗa iyalai biyu masu arziki na ƙasar.
Ioannes, \VUgras{angon} \textsuperscript{(18)}, \VUgras{ɗan} babban mai
gona ne, yanzu kuma ya zama \VUgras{surukin} wani gwanin masana’anta.

%%K t04-c1-10-1 p
10. --- \VUgras{Ma’aurata} matasa ne, duk da haka tun da daɗewa aka
yi musu \VUgras{alkawari.} \VUgras{Amarya} \textsuperscript{(19)},
Martha, tana da kwarjini cikin \VUgras{farar rigarta} da aka yi wa ado
da furannin lemu.

%%K t04-c1-11-1 p
11. --- \VUgras{Surukinta} \textsuperscript{(20)} yana farin ciki
ƙwarai da ya sami \VUgras{matar ɗa} mai kirki haka, \VUgras{surukuwar}
\textsuperscript{(21)} Ioannes kuwa tana murmushi tana ganin
\VUgras{ƴarta} da kyau haka.

%%K t04-c1-12-1 p
12. --- A ƙarshen tebur, \VUgras{baƙi} \textsuperscript{(22)} sauran
suna zaune, wato \VUgras{dangi, abokai, ƴan baffa, ƴan goggo.}
\VUgras{Shugaban masu hidima} \textsuperscript{(23)}, da tawul a
hammata, yana yi musu hidima da himma.

%%K t04-tit-4 sub
\VUpk{10.2pt}{40}{Abokai.}
\VUsaut{3.0mm}

%%K t04-c1-13-1 p
13. --- A hannun dama na tebur, ga \VUgras{budurwa}
\textsuperscript{(24)}, \VUgras{abokiyar} amaryar, sai ɗan uwanta,
wanda yake \VUgras{abokin girmamawarta} \textsuperscript{(25)}. Yana
hira da \VUgras{budurwar girmamawarsa} \textsuperscript{(26)}, ƙanwar
ango, wadda ta zama \VUgras{ƴar surukar} Martha ta wurin wannan aure.
Budurwa ce mai kwarjini. Tana da kyawawan \VUgras{kayan ado,}
\VUgras{sarƙar lu’ulu’u,} da \VUgras{munduwar zinare.}

%%K t04-c1-14-1 p
14. --- Kusa da su, malamin da yake tsaye yana ɗaga hidorinsa
\VUgras{abokin} \textsuperscript{(27)} iyalin ne. Yana ɗaya daga cikin
\VUgras{shaidun ango.} Yana \VUgras{gaisuwar sha,} yana sha domin
\VUgras{lafiyar} ma’auratan yana kuma yi musu fatan farin ciki mai
dorewa. Sanye yake da kayan biki, da kwat mai wutsiya da
\VUgras{takalma masu ƙyalli} \textsuperscript{(28)}. Shi ne mai raka
\VUgras{kaka} \textsuperscript{(29)}, wadda take \VUgras{gwauruwa.}

%%K t04-c1-15-1 p
15. --- Saurayin da yake sanye da \VUgras{kwat mai wutsiya}
\textsuperscript{(31)}, mai ɗan baƙin gashin baki, shi ne
\VUgras{surukin} \textsuperscript{(30)} Ioannes. Injiniya ne fitacce.
Shi ne maƙwabcin wata \VUgras{matashiyar mace} \textsuperscript{(33)}
mai kwarjini ƙwarai, wato \VUgras{yayar} Martha kuma mahaifiyar
ƙaramar yarinyar nan mai kyau wadda take zuwa, hannu da hannu da ɗan
baffanta, ta miƙa fure ta kuma \VUgras{yi masa fatan barka da
yamma.} Wannan mace tana da kyawawan \VUgras{ƴan kunne} da
\VUgras{ɗaurin kai} mai kwarjini.

%%K t04-c1-16-1 p
16. --- Ƙaramin ɗan baffan nan, mai ban mamaki saboda \VUgras{rigar
samansa} \textsuperscript{(36)} da \VUgras{gajeren wandonsa}
\textsuperscript{(37)} mai kumbura, abin tausayi ne ƙwarai.
\VUgras{Maraya} \textsuperscript{(35)} ne. Tun yana ƙarami ya rasa
mahaifinsa da mahaifiyarsa. Baffansa shi ne \VUgras{mai riƙonsa}
\textsuperscript{(38)}, wanda bai yi aure ba domin ya reni yaron nan
mai tausayi. Mutum ne mai kirki. Ɗan sanƙo ne kuma ba ya gani da
nisa ; yana amfani da \VUgras{tabarau mai maƙale.}

%%K t04-tit-5 sub
\VUsaut{2.4mm}
\VUpk{10.2pt}{40}{Kayan zaki.}
\VUsaut{3.0mm}

%%K t04-c1-17-1 p
17. --- A bayan masu biki, a kan tebur da aka rufe da \VUgras{zanen
tebur,} an shirya \VUgras{kayan zaki} \textsuperscript{(39)}. A cikin
babbar kwarya, ga ƴaƴan itace, \VUgras{fis} \textsuperscript{(40)},
\VUgras{fir} \textsuperscript{(41)}, \VUgras{tuffa}
\textsuperscript{(42)}, \VUgras{abirkot} \textsuperscript{(43)} da
\VUgras{inabi} \textsuperscript{(44)}. A kusa, \VUgras{abarba}
\textsuperscript{(45)}, \VUgras{kek} \textsuperscript{(46)} mai kyau,
\VUgras{kwalaben giya} \textsuperscript{(48)} da \VUgras{ruwan inabi mai
kumfa} \textsuperscript{(49)} da kuma jerin \VUgras{faranti}
\textsuperscript{(47)} masu tsabta.

%%K t04-c1-18-1 p
18. --- \VUgras{Mai riƙe da kwalabe} \textsuperscript{(50)} yana buɗe
\VUgras{kwalaban} \textsuperscript{(52)} tsohuwar giya da \VUgras{abin
cire matoshi} \textsuperscript{(51)}. \VUgras{Matoshin}
\textsuperscript{(53)} ya yi kama da mai wuyar cirewa ƙwarai. Wannan
mutum yana da \VUgras{tawul} \textsuperscript{(54)} a hammata.

%%K t04-c1-19-1 p
19. --- Bayan cin abincin dare, maza za su tafi ɗakin shan taba, mata
kuwa za su tafi su yi shirin rawa.

%%K t04-tit-6 sub
\VUsaut{5.0mm}
\VUcentre{11.4pt}{0}{\textit{Fage na Biyu.}}
\VUsaut{1.6mm}
\VUpk{11.4pt}{40}{Ranar haihuwar kaka.}
\VUsaut{2.6mm}

%%K t04-tit-7 sub
\VUpk{10.2pt}{40}{Ziyara.}
\VUsaut{3.0mm}

%%K t04-c2-01-1 p
1. --- Shekaru goma sun shuɗe, iyayen Ioannes sun tsufa suna kuma son
jikokinsu uku ƙwarai, wato \VUgras{jikoki maza} \textsuperscript{(2)}
biyu da \VUgras{jikanya} \textsuperscript{(3)} guda. Domin yau
\VUgras{ranar haihuwar kaka} ce, ƴaƴansu sun zo su yi masa fatan
barka da ranar haihuwa.

%%K t04-c2-02-1 p
2. --- Ƙaramin Petrus ƙarami ne har yanzu, shekarunsa uku kacal. Yana
ganin \VUgras{kyanwa} \textsuperscript{(1)} tana kusatowa. Yana son ya
shafa mata kyakkyawan \VUgras{gashinta} mai laushi da doguwar
\VUgras{wutsiyarta,} bai lura ba cewa a ƙarƙashin \VUgras{ƙafafunta}
masu kwarjini akwai farata masu kaifi ɓoye waɗanda wataƙila za su
kakkaɓe shi.

%%K t04-c2-03-1 p
3. --- Ƙanwarsa Gabriela, wadda take riƙe da hannunsa, ta fi shi
natsuwa ƙwarai, Marcellus kuwa da babbar \VUgras{alkyabbarsa}
\textsuperscript{(4)} da \VUgras{ɗamararsa} \textsuperscript{(5)} ta fata
ya yi kama da wanda ya girma, duk da gajeren wandonsa wanda yake
bayyana \VUgras{safarsa} \textsuperscript{(6)}.

%%K t04-c2-04-1 p
4. --- Mahaifinsu ya sa \VUgras{kwat na sanyi} \textsuperscript{(7)}
mai kauri na hunturu tare da \VUgras{kwalar gashi}
\textsuperscript{(8)}, a cikin \VUgras{aljihunsa} \textsuperscript{(9)}
kuma yana da ɗankwali. Matarsa tana da hularta ta gashi da aka yi wa
ado da \VUgras{gashin tsuntsu} \textsuperscript{(10)}, da
\VUgras{jaketarta} \textsuperscript{(11)}, da \VUgras{mayafin
gashinta} \textsuperscript{(12)} da \VUgras{abin ɗumama hannunta}
\textsuperscript{(13)}.

%%K t04-c2-05-1 p
5. --- Ana sanyi ƙwarai, domin hunturu ne. \VUgras{Dusar ƙanƙara}
\textsuperscript{(19)} ta yi ta sauka tsawon yini, rufin gidaje kuwa
duk ya yi fari. Tare da \VUgras{dare,} sanyi yana ƙara tsanani. A
sararin sama, \VUgras{wata} \textsuperscript{(17)} da \VUgras{taurari}
\textsuperscript{(18)} suna haske suna ratsa \VUgras{duhu.}

%%K t04-tit-8 sub
\VUsaut{4.2mm}
\VUpk{10.2pt}{40}{Ciwo.}
\VUsaut{3.0mm}

%%K t04-c2-06-1 p
6. --- \VUgras{Makaho} \textsuperscript{(20)} mai tausayin nan da yake
ratsa filin gaban \VUgras{kantin magani} \textsuperscript{(16)},
\VUgras{ƙaramin karensa} \textsuperscript{(21)} yana jagorantarsa,
bai da sa’a ko kaɗan. Iustinia, \VUgras{mai aikin gida}
\textsuperscript{(14)}, tana kawo daga ɗakin girki \VUgras{kwano}
\textsuperscript{(15)} na \VUgras{shayin ganye} mai zafi ƙwarai da aka
sa masa sukari da yawa.

%%K t04-c2-07-1 p
7. --- \VUgras{Kaka} \textsuperscript{(23)} tana son neman
\VUgras{tabarau mai hannunta} \textsuperscript{(26)} a kan tebur,
domin ta karanta \VUgras{takardar magani} \textsuperscript{(27)} da
\VUgras{likita} \textsuperscript{(25)} ya rubuta yanzu, sai ta tashi
daga babbar kujerarta, tana jingina a kan \VUgras{sandarta}
\textsuperscript{(24)}.

%%K t04-c2-08-1 p
8. --- Sau da yawa tana fama da \VUgras{ƙananan ciwo, jiri, mura,
ciwon haƙori,} ƙafafunta sun raunana idanunta kuma ba su da kyau
kuma. Duk da haka, tana yi wa tsohon \VUgras{likitanta} ba’a da so,
shi da yake so ya rubuta mata \VUgras{maganin sha}
\textsuperscript{(28)} kullum. Yau tana farin ciki ƙwarai domin tana da
ƙaramin iyalinta kewaye da ita.

%%K t04-c2-09-1 p
9. --- Ana jin daɗi cikin ɗakin da aka rufe da kyau. A cikin
\VUgras{murhun bango} \textsuperscript{(29)} na dutsen marmara,
\VUgras{wuta mai kyau} \textsuperscript{(30)} tana ci, mutum kuwa yana
jin daɗi cikin \VUgras{hasken} \VUgras{fitila} \textsuperscript{(31)}
mai sanyi da aka kunna yanzu.

%%K t04-tit-9 sub
\VUsaut{4.2mm}
\VUpk{10.2pt}{40}{Kaka.}
\VUsaut{3.0mm}

%%K t04-c2-10-1 p
10. --- Har \VUgras{kaka} \textsuperscript{(33)} da kansa ya manta cewa
yana da ciwo, cewa \VUgras{amosaninsa} ya ɗaure shi a kan
\VUgras{babbar kujerarsa} \textsuperscript{(34)} kuma cewa jiya ma
yana da \VUgras{zazzaɓi da suma.} Bai kuma tuna cewa a hunturun da ya
wuce ya sha \VUgras{ciwon huhu} kuma \VUgras{tari} mai ƙarfi da wahala
ya sa shi shan azaba ƙwarai.

%%K t04-c2-10-2 p
Duk da haka da wuya sosai aka warkar da shi. Yanzu yana cikin sauƙi.
Amma ƙafarsa wadda yake jingina a kan \VUgras{matashi}
\textsuperscript{(38)} da aka ajiye a kan ƙaramar \VUgras{kujera}
\textsuperscript{(39)} tana ba shi ciwo ita ma.

%%K t04-c2-11-1 p
11. --- Idan an lulluɓe shi da kula cikin \VUgras{rigar ɗakinsa}
\textsuperscript{(35)} mai ɗumi tare da \VUgras{maɓalli}
\textsuperscript{(36)} masu kauri na lu’ulu’u, kuma idan yana da
kusa da shi, domin ta karanta masa jarida, ƙaramar \VUgras{marayar}
\textsuperscript{(40)} nan sanye da farar \VUgras{hula,} da shuɗiyar
\VUgras{riga} \textsuperscript{(42)} da \VUgras{mayafi}
\textsuperscript{(41)}, ba ya gunaguni.

%%K t04-c2-12-1 p
12. --- Kowace shekara, idan \VUgras{kalanda} \textsuperscript{(43)}
ya sake nuna \VUgras{kwanan watan} ranar farko ga \VUgras{Disamba,}
yakan jira \VUgras{jikokinsa} ba tare da haƙuri ba, domin ya san cewa
ba za su manta da wannan \VUgras{kwanan wata} mai muhimmanci ba.

%%K t04-c2-13-1 p
13. --- Yakan tuna da tausayi lokacin nan na dā sosai da shi da kansa
ya je ya yi wa \VUgras{babban hafsan} sarki mai ɗaukaka fatan barka da
biki, wanda \VUgras{hotonsa} \textsuperscript{(44)} yake rataye a
bango, kuma wanda yake \VUgras{kakan kakan} ƴaƴansa.

\VUsaut{6.0mm}
\VUfilet{22mm}
